\documentclass[10pt,a4paper]{amsart}
\usepackage{amsmath,amssymb,amsthm,graphicx}
\usepackage{verbatim}
\newenvironment{reference}{\comment}{\endcomment}
%\newenvironment{reference}{}{}
\newenvironment{slogan}{\comment}{\endcomment}
\newenvironment{history}{\comment}{\endcomment}

% For commutative diagrams we use Xy-pic
\usepackage[all]{xy}

% We use 2cell for 2-commutative diagrams.
\xyoption{2cell}
% Theorem environments.
%
\theoremstyle{plain}
\newtheorem{theorem}[subsection]{Theorem}
\newtheorem{proposition}[subsection]{Proposition}
\newtheorem{lemma}[subsection]{Lemma}

\theoremstyle{definition}
\newtheorem{definition}[subsection]{Definition}
\newtheorem{example}[subsection]{Example}
\newtheorem{exercise}[subsection]{Exercise}
\newtheorem{situation}[subsection]{Situation}

\theoremstyle{remark}
\newtheorem{remark}[subsection]{Remark}
\newtheorem{remarks}[subsection]{Remarks}

\numberwithin{equation}{subsection}

% Macros
%
\def\lim{\mathop{\mathrm{lim}}\nolimits}
\def\colim{\mathop{\mathrm{colim}}\nolimits}
\def\Spec{\mathop{\mathrm{Spec}}}
\def\Hom{\mathop{\mathrm{Hom}}\nolimits}
\def\Ext{\mathop{\mathrm{Ext}}\nolimits}
\def\SheafHom{\mathop{\mathcal{H}\!\mathit{om}}\nolimits}
\def\SheafExt{\mathop{\mathcal{E}\!\mathit{xt}}\nolimits}
\def\Sch{\mathit{Sch}}
\def\Mor{\mathop{\mathrm{Mor}}\nolimits}
\def\Ob{\mathop{\mathrm{Ob}}\nolimits}
\def\Sh{\mathop{\mathit{Sh}}\nolimits}
\def\NL{\mathop{N\!L}\nolimits}
\def\CH{\mathop{\mathrm{CH}}\nolimits}
\def\proetale{{pro\text{-}\acute{e}tale}}
\def\etale{{\acute{e}tale}}
\def\QCoh{\mathit{QCoh}}
\def\Ker{\mathop{\mathrm{Ker}}}
\def\Im{\mathop{\mathrm{Im}}}
\def\Coker{\mathop{\mathrm{Coker}}}
\def\Coim{\mathop{\mathrm{Coim}}}

% Boxtimes
%
\DeclareMathSymbol{\boxtimes}{\mathbin}{AMSa}{"02}

%
% Macros for moduli stacks/spaces
%
\def\QCohstack{\mathcal{QC}\!\mathit{oh}}
\def\Cohstack{\mathcal{C}\!\mathit{oh}}
\def\Spacesstack{\mathcal{S}\!\mathit{paces}}
\def\Quotfunctor{\mathrm{Quot}}
\def\Hilbfunctor{\mathrm{Hilb}}
\def\Curvesstack{\mathcal{C}\!\mathit{urves}}
\def\Polarizedstack{\mathcal{P}\!\mathit{olarized}}
\def\Complexesstack{\mathcal{C}\!\mathit{omplexes}}
% \Pic is the operator that assigns to X its picard group, usage \Pic(X)
% \Picardstack_{X/B} denotes the Picard stack of X over B
% \Picardfunctor_{X/B} denotes the Picard functor of X over B
\def\Pic{\mathop{\mathrm{Pic}}\nolimits}
\def\Picardstack{\mathcal{P}\!\mathit{ic}}
\def\Picardfunctor{\mathrm{Pic}}
\def\Deformationcategory{\mathcal{D}\!\mathit{ef}}
\usepackage{cleveref}
\crefname{originallemma}{original lemma}{original lemmas}
\Crefname{originallemma}{Original lemma}{Original lemmas}
\makeatletter
\newlabel{algebra-lemma-descent-regular}{{07NG}{source}{}{sourcepointer-1}{}}
\newlabel{algebra-lemma-descent-regular@cref}{{[originallemma][0][]07NG}{[1][0][]source}{}{}{}}
\makeatother
\crefname{originallemma}{original lemma}{original lemmas}
\Crefname{originallemma}{Original lemma}{Original lemmas}
\makeatletter
\newlabel{coherent-lemma-proper-over-affine-cohomology-finite}{{02O6}{source}{}{sourcepointer-2}{}}
\newlabel{coherent-lemma-proper-over-affine-cohomology-finite@cref}{{[originallemma][0][]02O6}{[1][0][]source}{}{}{}}
\makeatother
\crefname{originaldefinition}{original definition}{original definitions}
\Crefname{originaldefinition}{Original definition}{Original definitions}
\makeatletter
\newlabel{morphisms-definition-alteration}{{0AB0}{source}{}{sourcepointer-3}{}}
\newlabel{morphisms-definition-alteration@cref}{{[originaldefinition][0][]0AB0}{[1][0][]source}{}{}{}}
\makeatother
\crefname{originalproposition}{original proposition}{original propositions}
\Crefname{originalproposition}{Original proposition}{Original propositions}
\makeatletter
\newlabel{morphisms-proposition-generic-flatness}{{052A}{source}{}{sourcepointer-4}{}}
\newlabel{morphisms-proposition-generic-flatness@cref}{{[originalproposition][0][]052A}{[1][0][]source}{}{}{}}
\makeatother
\crefname{originallemma}{original lemma}{original lemmas}
\Crefname{originallemma}{Original lemma}{Original lemmas}
\makeatletter
\newlabel{properties-lemma-normal-integral-sections}{{0358}{source}{}{sourcepointer-5}{}}
\newlabel{properties-lemma-normal-integral-sections@cref}{{[originallemma][0][]0358}{[1][0][]source}{}{}{}}
\makeatother
\crefname{originallemma}{original lemma}{original lemmas}
\Crefname{originallemma}{Original lemma}{Original lemmas}
\makeatletter
\newlabel{properties-lemma-regular-normal}{{0569}{source}{}{sourcepointer-6}{}}
\newlabel{properties-lemma-regular-normal@cref}{{[originallemma][0][]0569}{[1][0][]source}{}{}{}}
\makeatother
\title{Finiteness of the normalization and a regular dense open from the existence of a regular alteration}
\author{The Stacks Project Authors}
\date{stacks-project:resolve.tex:local-snapshot-20261012; GNU FDL 1.2 or later}
\begin{document}
\maketitle
\noindent\textit{Editorial note.} All delivered mathematical text is the author's own, copied from the pinned Stacks Project chapter file resolve.tex; no statement or proof step was rewritten, reordered or removed.\par\medskip
\noindent\textit{Editorial note.} The counted proof is the author's complete argument: the finiteness of the normalization through the ring of global sections of the alteration, and the regular dense open obtained by descending regularity along generic flatness.\par\medskip
\noindent\textit{Editorial note.} The label lemma-regular-alteration-implies names the counted statement, and the proof environment immediately following it in the source is delivered verbatim as its proof.\par\medskip
\noindent\textit{Editorial note.} The counted statement is carried by the author's own lemma environment in the chapter source.\par\medskip
\noindent\textit{Editorial note.} Pre-open gate: python operations/agent-authoring/tools/precard\_gate.py --source resolve.tex --label lemma-regular-alteration-implies --start 3746 --end 3773 --local-id 28107 returned RESULT=OK (L28107 free, occupancy FREE, public interval clear). Pre-assembly tag self-check: canonical tag resolved and not present in reference/public-items.jsonl.\par\medskip
\noindent\textit{Editorial note.} The statement and proof are copied from the pinned chapter file; the cited inputs are kept as external pointers with their canonical Stacks tags.\par\medskip
\noindent\textit{Allowed standard prerequisites (editorial).}\begin{itemize}
\item Finiteness of the global sections of a proper scheme over an affine base. Cohomology of Schemes, the lemma that the cohomology of a proper scheme over an affine base is finite over the base. Shows that the ring of global sections of the alteration is a finite module over the base ring.
\item Global sections of a normal integral scheme as a normal domain. Properties, the lemmas that a regular scheme is normal and that the global sections of a normal integral scheme form a normal domain. Places the normalization of the base inside the ring of global sections of the alteration as a subring.
\item Finiteness of the morphism over a nonempty open and generic flatness. Morphisms, the definition of an alteration and the generic flatness result. Produces an open over which the alteration is finite and flat.
\item Descent of regularity along a faithfully flat finite morphism. Algebra, the lemma descending regularity along a faithfully flat map. Concludes regularity of the open subscheme from regularity of its preimage.
\end{itemize}\medskip
\ifdefined\hypertarget\hypertarget{sourcepointer-1}{}\fi
\noindent\textit{Source reference.} Original lemma 07NG: sources/stacks-chow/tags.txt, source lines 9761--9761. Prerequisite pointer invoked at proof line(s) 3772; it names a standard fact of a sibling chapter that the counted argument uses, and every step of the counted proof built on it is retained verbatim in the delivered body.\par\medskip
\ifdefined\hypertarget\hypertarget{sourcepointer-2}{}\fi
\noindent\textit{Source reference.} Original lemma 02O6: sources/stacks-chow/tags.txt, source lines 3505--3505. Prerequisite pointer invoked at proof line(s) 3758; it names a standard fact of a sibling chapter that the counted argument uses, and every step of the counted proof built on it is retained verbatim in the delivered body.\par\medskip
\ifdefined\hypertarget\hypertarget{sourcepointer-3}{}\fi
\noindent\textit{Source reference.} Original definition 0AB0: sources/stacks-chow/tags.txt, source lines 13003--13003. Prerequisite pointer invoked at proof line(s) 3768; it names a standard fact of a sibling chapter that the counted argument uses, and every step of the counted proof built on it is retained verbatim in the delivered body.\par\medskip
\ifdefined\hypertarget\hypertarget{sourcepointer-4}{}\fi
\noindent\textit{Source reference.} Original proposition 052A: sources/stacks-chow/tags.txt, source lines 6563--6563. Prerequisite pointer invoked at proof line(s) 3770; it names a standard fact of a sibling chapter that the counted argument uses, and every step of the counted proof built on it is retained verbatim in the delivered body.\par\medskip
\ifdefined\hypertarget\hypertarget{sourcepointer-5}{}\fi
\noindent\textit{Source reference.} Original lemma 0358: sources/stacks-chow/tags.txt, source lines 4119--4119. Prerequisite pointer invoked at proof line(s) 3762; it names a standard fact of a sibling chapter that the counted argument uses, and every step of the counted proof built on it is retained verbatim in the delivered body.\par\medskip
\ifdefined\hypertarget\hypertarget{sourcepointer-6}{}\fi
\noindent\textit{Source reference.} Original lemma 0569: sources/stacks-chow/tags.txt, source lines 6702--6702. Prerequisite pointer invoked at proof line(s) 3760; it names a standard fact of a sibling chapter that the counted argument uses, and every step of the counted proof built on it is retained verbatim in the delivered body.\par\medskip
\begin{lemma}
\label{lemma-regular-alteration-implies}
Let $Y$ be a Noetherian integral scheme. Assume there exists an alteration
$f : X \to Y$ with $X$ regular. Then the normalization $Y^\nu \to Y$
is finite and $Y$ has a dense open which is regular.
\end{lemma}

\begin{proof}
It suffices to prove this when $Y = \Spec(A)$ where $A$ is a Noetherian domain.
Let $B$ be the integral closure of $A$ in its fraction field.
Set $C = \Gamma(X, \mathcal{O}_X)$. By
Cohomology of Schemes, Lemma
\ref{coherent-lemma-proper-over-affine-cohomology-finite}
we see that $C$ is a finite $A$-module. As $X$ is normal
(Properties, Lemma \ref{properties-lemma-regular-normal})
we see that $C$ is normal domain
(Properties, Lemma \ref{properties-lemma-normal-integral-sections}).
Thus $B \subset C$ and we conclude that $B$ is finite over $A$
as $A$ is Noetherian.

\medskip\noindent
There exists a nonempty open $V \subset Y$ such that $f^{-1}V \to V$
is finite, see Morphisms, Definition \ref{morphisms-definition-alteration}.
After shrinking $V$ we may assume that $f^{-1}V \to V$ is flat
(Morphisms, Proposition \ref{morphisms-proposition-generic-flatness}).
Thus $f^{-1}V \to V$ is faithfully flat. Then $V$ is regular by
Algebra, Lemma \ref{algebra-lemma-descent-regular}.
\end{proof}

\end{document}
